\subsection{Why curves of second degree appear}

The previous chapter used first-degree equations to describe lines and planes. A
new phenomenon appears when the defining condition contains Euclidean distance:
squares enter automatically.

\begin{corebox}[The question of this chapter]
Can one choose coordinates so that the geometry of circles, ellipses, parabolas,
and hyperbolas becomes visible rather than hidden inside algebra?
\end{corebox}

We begin with the familiar Cartesian pictures. They provide a reference point,
not the final language. The real goal is to place the origin at a focus and ask
for the distance to the curve in each direction.

\begin{center}
\begin{tikzpicture}[scale=0.72,>=stealth]
    \begin{scope}[xshift=-5.4cm]
        \draw[axis,->] (-2.4,0)--(2.6,0);
        \draw[axis,->] (0,-2.2)--(0,2.4);
        \draw[subset] (0,0) circle (1.5);
        \node at (0,-2.7) {circle};
    \end{scope}
    \begin{scope}[xshift=-1.8cm]
        \draw[axis,->] (-2.5,0)--(2.7,0);
        \draw[axis,->] (0,-2.2)--(0,2.4);
        \draw[subset] (0,0) ellipse (2.0 and 1.2);
        \node at (0,-2.7) {ellipse};
    \end{scope}
    \begin{scope}[xshift=2.1cm]
        \draw[axis,->] (-1.5,0)--(3.1,0);
        \draw[axis,->] (0,-2.2)--(0,2.4);
        \draw[subset,domain=-2:2,samples=100] plot ({0.45*\x*\x},{\x});
        \node at (0.8,-2.7) {parabola};
    \end{scope}
    \begin{scope}[xshift=6.0cm]
        \draw[axis,->] (-2.8,0)--(3.0,0);
        \draw[axis,->] (0,-2.2)--(0,2.4);
        \draw[subset,domain=1.05:2.65,samples=80] plot (\x,{0.75*sqrt(\x*\x-1)});
        \draw[subset,domain=1.05:2.65,samples=80] plot (\x,{-0.75*sqrt(\x*\x-1)});
        \draw[subset,domain=-2.65:-1.05,samples=80] plot (\x,{0.75*sqrt(\x*\x-1)});
        \draw[subset,domain=-2.65:-1.05,samples=80] plot (\x,{-0.75*sqrt(\x*\x-1)});
        \node at (0,-2.7) {hyperbola};
    \end{scope}
\end{tikzpicture}
\end{center}

\begin{interpretationbox}[A map of the argument]
First we identify the four shapes in Cartesian coordinates. Then we construct
polar coordinates. Finally one focus--directrix condition produces one polar
formula whose parameter decides which conic appears.
\end{interpretationbox}
